\documentclass[10pt,twocolumn,letterpaper]{article}

\usepackage{cvpr}
\usepackage{times}
\usepackage{epsfig}
\usepackage{graphicx}
\usepackage{amsmath}
\usepackage{amssymb}
\usepackage{booktabs}
\usepackage{multirow}
\usepackage[pagebackref]{hyperref}
\usepackage{listings}
\usepackage{xcolor}
\usepackage{longtable}

% Code listing style
\lstset{
    basicstyle=\ttfamily\small,
    breaklines=true,
    frame=single,
    backgroundcolor=\color{gray!10}
}

\def\cvprPaperID{7}
\def\paperID{7}
\def\confName{CVPR}
\def\confYear{2026}

\begin{document}

\title{Supplementary Materials: Emotional Vocabulary as Semantic Grounding}

\author{Scott Boudreaux\\
Elyan Labs\\
{\tt\small scott@elyanlabs.com}
}

\maketitle

This document provides supplementary materials for ``Emotional Vocabulary as Semantic Grounding: How Language Register Affects Diffusion Efficiency in Video Generation.'' We include: (A)~full per-seed results for all 35 matched pairs, (B)~complete STOCK/NEURO prompt pairs, (C)~the Prompt Translator motion-to-emotion dictionary, (D)~statistical test details, and (E)~cross-model experimental parameters.

%----------------------------------------------------------------------
\section{Full Per-Seed Results}
\label{sec:perseed}
%----------------------------------------------------------------------

Table~\ref{tab:full_results} reports file sizes (bytes), size delta~(\%), and mean frame-level LPIPS for all 35 matched STOCK/NEURO pairs. Seeds are drawn from the range $42{,}424{,}242$--$42{,}428{,}242$ in increments of $1{,}000$. Each LPIPS value is the mean across 24 frames per video pair.

\begin{table*}[t]
\centering
\small
\caption{Complete per-seed results for all 35 matched pairs. STOCK uses 30~steps / guidance~7.5; NEURO uses 24~steps / guidance~8.0. Size $\Delta$ is computed as $(\text{NEURO} - \text{STOCK})/\text{STOCK} \times 100$.}
\label{tab:full_results}
\begin{tabular}{llrrrc}
\toprule
\textbf{Emotional Arc} & \textbf{Seed} & \textbf{STOCK (bytes)} & \textbf{NEURO (bytes)} & \textbf{Size $\Delta$ (\%)} & \textbf{LPIPS} \\
\midrule
\multirow{5}{*}{debate\_passion}
 & 42424242 & 564,980 & 508,390 & $-$10.0 & 0.600 \\
 & 42425242 & 676,678 & 405,148 & $-$40.1 & 0.511 \\
 & 42426242 & 675,138 & 661,622 & $-$2.0 & 0.504 \\
 & 42427242 & 701,512 & 566,978 & $-$19.2 & 0.557 \\
 & 42428242 & 665,464 & 302,488 & $-$54.5 & 0.552 \\
\cmidrule{2-6}
 & \textit{Mean} & \textit{656,954} & \textit{488,925} & \textit{$-$25.6} & \textit{0.545} \\
\midrule
\multirow{5}{*}{debate\_tension}
 & 42424242 & 560,590 & 572,210 & $+$2.1 & 0.010 \\
 & 42425242 & 640,972 & 672,998 & $+$5.0 & 0.021 \\
 & 42426242 & 640,914 & 836,696 & $+$30.5 & 0.079 \\
 & 42427242 & 659,404 & 664,694 & $+$0.8 & 0.017 \\
 & 42428242 & 661,228 & 639,612 & $-$3.3 & 0.019 \\
\cmidrule{2-6}
 & \textit{Mean} & \textit{632,622} & \textit{677,242} & \textit{$+$7.1} & \textit{0.029} \\
\midrule
\multirow{5}{*}{elyan\_claude\_focus}
 & 42424242 & 860,000 & 911,312 & $+$6.0 & 0.025 \\
 & 42425242 & 860,712 & 339,426 & $-$60.6 & 0.597 \\
 & 42426242 & 949,840 & 510,758 & $-$46.2 & 0.559 \\
 & 42427242 & 983,726 & 634,312 & $-$35.5 & 0.523 \\
 & 42428242 & 821,832 & 630,912 & $-$23.2 & 0.623 \\
\cmidrule{2-6}
 & \textit{Mean} & \textit{895,222} & \textit{605,344} & \textit{$-$32.4} & \textit{0.465} \\
\midrule
\multirow{5}{*}{elyan\_sophia\_focus}
 & 42424242 & 868,184 & 1,012,660 & $+$16.6 & 0.045 \\
 & 42425242 & 949,600 & 1,001,102 & $+$5.4 & 0.467 \\
 & 42426242 & 998,280 & 1,076,892 & $+$7.9 & 0.584 \\
 & 42427242 & 1,109,272 & 1,150,352 & $+$3.7 & 0.558 \\
 & 42428242 & 851,244 & 1,040,128 & $+$22.2 & 0.576 \\
\cmidrule{2-6}
 & \textit{Mean} & \textit{955,316} & \textit{1,056,227} & \textit{$+$10.6} & \textit{0.446} \\
\midrule
\multirow{5}{*}{sophia\_contemplation}
 & 42424242 & 244,698 & 237,394 & $-$3.0 & 0.008 \\
 & 42425242 & 265,550 & 267,914 & $+$0.9 & 0.008 \\
 & 42426242 & 296,974 & 322,404 & $+$8.6 & 0.011 \\
 & 42427242 & 250,848 & 242,754 & $-$3.2 & 0.013 \\
 & 42428242 & 212,636 & 226,070 & $+$6.3 & 0.011 \\
\cmidrule{2-6}
 & \textit{Mean} & \textit{254,141} & \textit{259,307} & \textit{$+$2.0} & \textit{0.010} \\
\midrule
\multirow{5}{*}{sophia\_determination}
 & 42424242 & 222,346 & 215,214 & $-$3.2 & 0.011 \\
 & 42425242 & 299,428 & 303,092 & $+$1.2 & 0.004 \\
 & 42426242 & 303,892 & 300,658 & $-$1.1 & 0.008 \\
 & 42427242 & 239,642 & 248,470 & $+$3.7 & 0.012 \\
 & 42428242 & 220,598 & 214,018 & $-$3.0 & 0.006 \\
\cmidrule{2-6}
 & \textit{Mean} & \textit{257,181} & \textit{256,290} & \textit{$-$0.3} & \textit{0.008} \\
\midrule
\multirow{5}{*}{sophia\_realization}
 & 42424242 & 257,538 & 217,848 & $-$15.4 & 0.025 \\
 & 42425242 & 310,924 & 312,302 & $+$0.4 & 0.008 \\
 & 42426242 & 322,000 & 326,024 & $+$1.2 & 0.012 \\
 & 42427242 & 296,332 & 274,470 & $-$7.4 & 0.013 \\
 & 42428242 & 249,110 & 204,154 & $-$18.0 & 0.016 \\
\cmidrule{2-6}
 & \textit{Mean} & \textit{287,181} & \textit{266,960} & \textit{$-$7.0} & \textit{0.015} \\
\bottomrule
\end{tabular}
\end{table*}

\textbf{Observations.} File size delta shows high per-seed variance (e.g., debate\_passion ranges from $-$54.5\% to $-$2.0\%), which is expected given the stochastic nature of diffusion sampling. LPIPS, by contrast, is highly stable within each arc ($\sigma < 0.035$ for complex arcs, $\sigma < 0.006$ for solo arcs), confirming that perceptual equivalence or divergence is a robust property of the prompt condition rather than seed noise.

%----------------------------------------------------------------------
\section{Full Prompt Pairs}
\label{sec:prompts}
%----------------------------------------------------------------------

All seven STOCK/NEURO prompt pairs used in the benchmark are listed below. Each pair shares the same source image and seed set. STOCK prompts describe external motion; NEURO prompts describe internal emotional states.

\subsection{Arc 1: Realization (sophia\_realization)}

\textbf{STOCK:}
\begin{lstlisting}
Victorian woman portrait, subtle head
movement, slight smile, blinking eyes,
warm lighting
\end{lstlisting}

\textbf{NEURO:}
\begin{lstlisting}
The young woman's eyes brighten with
quiet realization, a knowing smile
forming as inspiration takes hold,
warmth spreading across her expression
\end{lstlisting}

\subsection{Arc 2: Contemplation (sophia\_contemplation)}

\textbf{STOCK:}
\begin{lstlisting}
Victorian woman portrait, looking
thoughtful, gentle movements, soft
lighting
\end{lstlisting}

\textbf{NEURO:}
\begin{lstlisting}
Her gaze turns inward with deep
contemplation, a subtle shift from
curiosity to understanding, quiet
wisdom settling in her features
\end{lstlisting}

\subsection{Arc 3: Determination (sophia\_determination)}

\textbf{STOCK:}
\begin{lstlisting}
Victorian woman portrait, serious
expression, focused look, slight
movement
\end{lstlisting}

\textbf{NEURO:}
\begin{lstlisting}
Quiet determination hardens in her
eyes, jaw setting with newfound
resolve, inner fire building behind
composed exterior
\end{lstlisting}

\subsection{Arc 4: Confidence (elyan\_sophia\_focus)}

\textbf{STOCK:}
\begin{lstlisting}
Victorian exhibition, woman working
on machine, man watching, gaslight
flickering
\end{lstlisting}

\textbf{NEURO:}
\begin{lstlisting}
The young woman works with fierce
concentration, confident hands moving
with purpose, quiet authority radiating
as she masters the brass machinery
\end{lstlisting}

\subsection{Arc 5: Respect (elyan\_claude\_focus)}

\textbf{STOCK:}
\begin{lstlisting}
Victorian exhibition, older man
gesturing, woman at machine, warm
lighting
\end{lstlisting}

\textbf{NEURO:}
\begin{lstlisting}
The older gentleman's skepticism
softens to grudging respect, pride
wounded but giving way to reluctant
admiration
\end{lstlisting}

\subsection{Arc 6: Passion (debate\_passion)}

\textbf{STOCK:}
\begin{lstlisting}
Two people in conversation, gesturing,
fireplace glowing, Victorian study
\end{lstlisting}

\textbf{NEURO:}
\begin{lstlisting}
Passionate intellectual exchange,
conviction burning in their eyes,
the electricity of clashing ideas
filling the air between them
\end{lstlisting}

\subsection{Arc 7: Tension (debate\_tension)}

\textbf{STOCK:}
\begin{lstlisting}
Two people talking, subtle movements,
warm firelight, period room
\end{lstlisting}

\textbf{NEURO:}
\begin{lstlisting}
Tension crackling between them,
unspoken challenge in their gazes,
the air thick with intellectual rivalry
\end{lstlisting}

%----------------------------------------------------------------------
\section{Prompt Translator Motion-to-Emotion Dictionary}
\label{sec:pt_dict}
%----------------------------------------------------------------------

Table~\ref{tab:pt_dictionary} presents the complete 22-entry motion-to-emotion dictionary used by the Prompt Translator. Each mapping replaces a literal motion descriptor with an emotionally-grounded equivalent that activates denser embedding regions. These mappings were derived empirically through iterative A/B testing.

\begin{table}[t]
\centering
\small
\caption{Prompt Translator motion-to-emotion dictionary. Each literal motion term (left) is replaced by its emotional equivalent (right) during prompt translation.}
\label{tab:pt_dictionary}
\begin{tabular}{p{2.8cm}p{4.2cm}}
\toprule
\textbf{Literal Motion} & \textbf{Emotional Equivalent} \\
\midrule
head movement & subtle shift in attention \\
head tilt & curiosity awakening \\
nods & quiet agreement settling \\
shakes head & gentle denial forming \\
blinks & moment of processing \\
slight smile & knowing warmth emerging \\
frowns & concern deepening \\
looks at camera & awareness sharpening \\
stares & intensity building \\
eye movement & attention drifting \\
hand movement & gesture carrying emotional weight \\
gestures & expression through motion \\
speaks & conviction forming in words \\
turns & attention shifting with purpose \\
leans forward & engagement intensifying \\
leans back & contemplative withdrawal \\
moves closer & warmth drawing near \\
moves away & reluctant distancing \\
raises eyebrows & surprise dawning \\
squints & scrutiny deepening \\
jaw clenches & resolve hardening \\
shoulders tense & burden settling \\
\bottomrule
\end{tabular}
\end{table}

\textbf{Design rationale.} Each mapping converts a biomechanical description (what the body \emph{does}) into a psychological description (what the character \emph{feels}). This shift exploits the observation that text encoders trained on web-scale data contain denser representations for emotional states than for motion mechanics, since natural language captions overwhelmingly describe the \emph{meaning} of expressions rather than their biomechanics.

%----------------------------------------------------------------------
\section{Statistical Test Details}
\label{sec:stats}
%----------------------------------------------------------------------

We report four primary statistical tests. All tests use standard significance thresholds ($\alpha = 0.05$).

\subsection{File Size: Wilcoxon Signed-Rank Test}

\textbf{Hypothesis.} $H_0$: The distribution of file size differences (NEURO $-$ STOCK) is symmetric about zero.

\textbf{Data.} $n = 35$ matched pairs. Differences range from $-362{,}530$ bytes to $+195{,}782$ bytes.

\textbf{Result.} $W = 264$, $p = 0.413$ (two-sided). We \textbf{fail to reject} $H_0$. The overall file size difference is not statistically significant, reflecting the high per-seed variance in compressed video output. This motivates our use of LPIPS as the primary quality metric rather than file size.

\textbf{Interpretation.} File size serves as a rough efficiency proxy but is too noisy for individual-pair significance. The aggregate $-$8.3\% mean reduction reflects a real trend, but the Wilcoxon test correctly identifies that the effect is not uniformly directional across all arcs (e.g., tension and confidence show increases).

\subsection{Solo Portrait LPIPS: One-Sample $t$-Test}

\textbf{Hypothesis.} $H_0$: Mean solo-portrait LPIPS $\geq 0.1$ (perceptual equivalence threshold from Zhang et al.~\cite{zhang2018}).

\textbf{Data.} \textit{Primary analysis}: $n = 15$ core solo-portrait pairs (arcs: contemplation $\times 5$, determination $\times 5$, realization $\times 5$). These arcs contain only single-subject scenes with unambiguous solo framing.

\textbf{Result (primary, $n = 15$).}
\begin{align}
\bar{x} &= 0.011, \quad s = 0.005 \nonumber \\
t &= \frac{0.011 - 0.1}{0.005 / \sqrt{15}} = -69.59 \nonumber \\
p &< 10^{-19} \quad \text{(one-sided)}
\end{align}

We \textbf{reject} $H_0$ with overwhelming significance. Solo-portrait NEURO outputs are perceptually equivalent to STOCK outputs, validating the design choice of allocating 20\% fewer diffusion steps to emotionally-grounded prompts.

\textit{Extended analysis}: Including the tension arc ($n{=}5$, LPIPS $= 0.029 \pm 0.025$) and two low-LPIPS seeds from complex arcs yields $n = 22$, $\bar{x} = 0.017$, $s = 0.016$, $t = -23.58$, $p < 10^{-8}$. However, the two complex-arc seeds were selected post-hoc based on low LPIPS values, introducing selection bias. We therefore report the $n = 15$ analysis as primary.

\subsection{Ablation: One-Sample $t$-Test}

\textbf{Hypothesis.} $H_0$: Mean ablation LPIPS $\geq 0.1$ when STOCK and NEURO use identical parameters (30~steps, guidance~7.5).

\textbf{Data.} $n = 9$ matched pairs across 3 arcs (respect, realization, determination) $\times$ 3 seeds.

\textbf{Result.}
\begin{align}
\bar{x} &= 0.068, \quad s = 0.029 \nonumber \\
t &= \frac{0.068 - 0.1}{0.029 / \sqrt{9}} = -3.14 \nonumber \\
p &= 0.014 \quad \text{(one-sided)}
\end{align}

We \textbf{reject} $H_0$. Even with identical computational budget, emotional prompts produce perceptually equivalent outputs. This confirms the effect is prompt-driven, not an artifact of the step reduction.

\subsection{Effect Size: Cohen's $d$}

To quantify the practical magnitude of the LPIPS difference from the 0.1 threshold:

\begin{equation}
d = \frac{|\bar{x} - \mu_0|}{s} = \frac{|0.011 - 0.1|}{0.005} = 17.8
\end{equation}

\noindent where $\mu_0 = 0.1$ is the perceptual equivalence threshold, computed over the $n = 15$ core solo-portrait pairs. A Cohen's $d$ of 17.8 indicates that solo-portrait LPIPS values lie approximately 18 standard deviations below the perceptual equivalence threshold---reflecting the fact that STOCK and NEURO outputs are nearly pixel-identical for single-subject scenes, not merely ``close enough.'' For reference, the extended $n = 22$ set yields $d = 5.19$, still an overwhelmingly large effect.

\subsection{Summary of Statistical Evidence}

\begin{table}[h]
\centering
\small
\caption{Summary of all statistical tests.}
\label{tab:stats_summary}
\begin{tabular}{@{}lccl@{}}
\toprule
\textbf{Test} & \textbf{Stat.} & \textbf{$p$} & \textbf{Result} \\
\midrule
File size (Wilcoxon) & $W\!=\!264$ & 0.413 & Fail to rej. \\
Solo LPIPS ($n\!=\!15$) & $t\!=\!{-}69.6$ & $<\!10^{-19}$ & Reject $H_0$ \\
Ablation ($n\!=\!9$) & $t\!=\!{-}3.14$ & 0.014 & Reject $H_0$ \\
Cohen's $d$ ($n\!=\!15$) & $d\!=\!17.8$ & --- & Very large \\
\bottomrule
\end{tabular}
\end{table}

%----------------------------------------------------------------------
\section{Cross-Model Experimental Parameters}
\label{sec:crossmodel}
%----------------------------------------------------------------------

We validated the emotional prompting effect on two additional architectures. Table~\ref{tab:crossmodel_params} reports the experimental configurations.

\subsection{AnimateDiff}

\begin{table}[h]
\centering
\small
\caption{AnimateDiff experimental parameters.}
\label{tab:animatediff_params}
\begin{tabular}{lcc}
\toprule
\textbf{Parameter} & \textbf{STOCK} & \textbf{NEURO} \\
\midrule
Model & \multicolumn{2}{c}{AnimateDiff v2 + SDXL} \\
Text encoder & \multicolumn{2}{c}{CLIP ViT-L/14 (400M)} \\
Steps & 30 & 24 \\
Guidance scale & 7.5 & 8.0 \\
Resolution & \multicolumn{2}{c}{512 $\times$ 512} \\
Frames & \multicolumn{2}{c}{16} \\
Motion module & \multicolumn{2}{c}{v2\_lora\_PanLeft} \\
Seeds per arc & \multicolumn{2}{c}{3} \\
Arcs tested & \multicolumn{2}{c}{realization, determination, respect} \\
Total renders & \multicolumn{2}{c}{18} \\
\bottomrule
\end{tabular}
\end{table}

\textbf{Key findings.} CLIP-based AnimateDiff shows an architecture-dependent effect: solo portraits exhibit \emph{reversed} results where NEURO produces +30\% (realization) and +81\% (determination) \emph{larger} files. For complex multi-character scenes (respect), both architectures agree on $\sim$33\% efficiency gain from emotional prompting. This architecture dependence is explained by CLIP's contrastive training objective (image-text matching) versus Gemma~3's language understanding objective.

\subsection{SVD XT (Stable Video Diffusion)}

\begin{table}[h]
\centering
\small
\caption{SVD XT experimental parameters.}
\label{tab:svd_params}
\begin{tabular}{lc}
\toprule
\textbf{Parameter} & \textbf{Value} \\
\midrule
Model & SVD XT 1.1 \\
Conditioning & Image-only (no text) \\
Steps & 25 \\
Motion bucket ID & 127 \\
Resolution & 1024 $\times$ 576 \\
Frames & 25 \\
Augmentation level & 0.0 \\
Seeds tested & 6 \\
\bottomrule
\end{tabular}
\end{table}

\textbf{Key findings.} SVD~XT uses image-only conditioning without text prompts. File size coefficient of variation across 6~seeds was $\sim$2.7\%, establishing a baseline for natural stochastic variance. This confirms the emotional prompting effect \textbf{requires text conditioning} and is not an artifact of diffusion sampling noise.

\subsection{Cross-Model Results Summary}

\begin{table}[h]
\centering
\small
\caption{Cross-model validation results. ``$\Delta$'' denotes NEURO file size change relative to STOCK.}
\label{tab:crossmodel_params}
\begin{tabular}{lccc}
\toprule
\textbf{Scenario} & \textbf{LTX-2} & \textbf{AnimateDiff} & \textbf{SVD XT} \\
 & \textbf{(Gemma 3)} & \textbf{(CLIP)} & \textbf{(None)} \\
\midrule
Solo portraits & $-$5\% & $+$55\% & N/A \\
Complex scenes & $-$36\% & $-$33\% & N/A \\
Seed variance & --- & --- & 2.7\% CV \\
\bottomrule
\end{tabular}
\end{table}

\textbf{Universal result.} Complex emotional scenes with multiple characters benefit $\sim$33\% from emotional prompting regardless of text encoder architecture. This is the most robust use case for the technique. Solo-scene benefits are architecture-dependent: language understanding encoders (Gemma~3) show consistent gains, while contrastive encoders (CLIP) show reversed effects on simple scenes.

%----------------------------------------------------------------------
\section{Generation Parameters}
\label{sec:gen_params}
%----------------------------------------------------------------------

For completeness, Table~\ref{tab:gen_params_full} reports all generation parameters for the primary LTX-2 benchmark.

\begin{table}[h]
\centering
\small
\caption{Full LTX-2 generation parameters.}
\label{tab:gen_params_full}
\begin{tabular}{lcc}
\toprule
\textbf{Parameter} & \textbf{STOCK} & \textbf{NEURO} \\
\midrule
Diffusion steps & 30 & 24 \\
Guidance scale & 7.5 & 8.0 \\
Max shift & 2.05 & 2.10 \\
Base shift & 0.95 & 0.98 \\
Resolution & 512 $\times$ 320 & 512 $\times$ 320 \\
Frames & 49 & 49 \\
Output format & WebP & WebP \\
Seeds per arc & 5 & 5 \\
Arcs & 7 & 7 \\
Total renders & 35 & 35 \\
\midrule
Model & \multicolumn{2}{c}{LTX-2 (19B params)} \\
Text encoder & \multicolumn{2}{c}{Gemma 3 12B} \\
Pipeline & \multicolumn{2}{c}{ComfyUI} \\
GPU & \multicolumn{2}{c}{V100 32GB} \\
LPIPS backbone & \multicolumn{2}{c}{AlexNet} \\
\bottomrule
\end{tabular}
\end{table}

%----------------------------------------------------------------------

\section{Independent Encoder Validation}

To test whether the embedding density pattern generalizes beyond Gemma 3, we computed pairwise cosine distances for emotional vs.\ literal vocabulary using \texttt{all-MiniLM-L6-v2} (384-dim, 22M parameters). Results:

\begin{table}[h]
\centering
\small
\begin{tabular}{lcc}
\toprule
\textbf{Vocab.} & \textbf{Mean Dist.} & \textbf{Encoder} \\
\midrule
Emotional & 0.713 & MiniLM-L6 \\
Literal & 0.722 & MiniLM-L6 \\
\midrule
Emotional & 0.225 & Gemma 3 12B \\
Literal & 0.269 & Gemma 3 12B \\
\bottomrule
\end{tabular}
\caption{Emotional vocabulary clusters tighter in both encoders, though the effect is much stronger in Gemma 3 (16\% vs.\ 1.2\%).}
\end{table}

The consistent direction across two very different encoders (22M vs 12B params) suggests the density pattern is a genuine property of emotional vocabulary in pretrained models, not an artifact of a specific architecture. The much stronger effect in Gemma 3 is consistent with our finding that large-scale language model encoders capture emotional semantics more deeply than smaller models.

\subsection{STOCK vs NEURO Prompt Similarity}

Cosine similarity between matched STOCK/NEURO prompt pairs averages 0.413 (range: 0.30--0.57), confirming the prompts are semantically distinct---not mere paraphrases. STOCK prompts cluster more tightly as a group (mean distance 0.431) due to shared template vocabulary (``portrait,'' ``lighting,'' ``Victorian''). NEURO prompts are more diverse at the sentence level (mean distance 0.605) because each encodes a unique emotional arc.

\section{Human Evaluation Protocol (Proposed)}

We acknowledge the absence of human evaluation as a limitation. Below we describe a protocol for future work, designed to be reproducible and statistically powered.

\subsection{Design}

\textbf{Task}: Two-alternative forced choice (2AFC). Evaluators view a STOCK/NEURO video pair (randomized left/right) and answer:

\begin{enumerate}
\item \textit{Quality}: ``Which video has higher visual quality?'' (or ``No difference'')
\item \textit{Expressiveness}: ``Which video conveys more emotion?'' (or ``No difference'')
\item \textit{Preference}: ``Which video do you prefer overall?'' (or ``No difference'')
\end{enumerate}

\textbf{Stimuli}: 14 pairs (7 arcs $\times$ 2 seeds), balanced across scene types. Videos displayed side-by-side at native resolution, looping.

\textbf{Participants}: Minimum 10 evaluators (5 with video/VFX expertise, 5 naive). Power analysis: with $n=10$ raters and 14 stimuli, a binomial test detects $\geq$70\% agreement at $\alpha = 0.05$ with power $> 0.80$.

\textbf{Analysis}: Per-question binomial test against chance (50\%). Inter-rater agreement via Fleiss' $\kappa$. Stratified by scene type (solo vs.\ complex).

\textbf{Hypothesis}: For solo portraits, we expect no significant quality preference (consistent with LPIPS $<$ 0.02). For complex scenes, we expect NEURO preference on expressiveness but not necessarily quality.

\section{CLIP Image-Text Alignment Scores}

As a complementary metric to LPIPS, we computed CLIP image-text similarity (ViT-B/32) between each video's first frame and its corresponding prompt text. Results reveal a dissociation between the two metrics:

Solo portraits: NEURO prompts achieve 13.5\% higher CLIP alignment (0.231 vs 0.204), suggesting emotional vocabulary provides better semantic grounding even by CLIP's contrastive criteria.

Complex scenes: STOCK prompts achieve 17.4\% higher CLIP alignment (0.296 vs 0.244), because their literal descriptions directly match visible scene content that CLIP can verify.

Overall mean: STOCK 0.257, NEURO 0.239. This confirms CLIP score and LPIPS measure different aspects of generation quality.

\begin{table}[h]
\centering
\small
\caption{Per-arc CLIP image-text alignment scores (ViT-B/32) between the first video frame and the corresponding prompt.}
\label{tab:clip_alignment}
\begin{tabular}{lcc}
\toprule
\textbf{Arc} & \textbf{STOCK} & \textbf{NEURO} \\
\midrule
sophia\_realization & 0.232 & \textbf{0.245} \\
sophia\_contemplation & 0.179 & \textbf{0.225} \\
sophia\_determination & 0.200 & \textbf{0.224} \\
elyan\_sophia\_focus & \textbf{0.241} & 0.226 \\
elyan\_claude\_focus & \textbf{0.268} & 0.212 \\
debate\_passion & \textbf{0.355} & 0.272 \\
debate\_tension & \textbf{0.320} & 0.268 \\
\bottomrule
\end{tabular}
\end{table}

\section{Diverse Image Generalization Details}
\label{app:diverse}

Section~5.8 of the main paper reports aggregate results for the diverse image benchmark. Here we provide the full prompt pairs and per-seed results.

\subsection{Prompt Pairs for Diverse Images}

Each image was rendered under both STOCK (literal motion descriptors) and NEURO (emotional vocabulary) conditions with identical seeds.

\begin{table*}[h]
\centering
\small
\begin{tabular}{p{2cm}p{6cm}p{6cm}}
\toprule
\textbf{Image} & \textbf{STOCK Prompt} & \textbf{NEURO Prompt} \\
\midrule
Mountain Lake (landscape) & Mountain lake scene, calm water with slight ripples, clouds moving slowly overhead, soft afternoon light, trees swaying gently & Serene tranquility washes over the mountain lake, nature holding its breath as golden light spills across still waters, a profound peace radiating from the ancient peaks \\
\midrule
Medieval Market (architecture) & Medieval marketplace, flags moving in wind, torch flames flickering, shadows shifting on cobblestones, warm evening light & The bustling energy of ancient commerce fills the marketplace, hopeful anticipation rippling through the crowd as torchlight dances with restless ambition across weathered stone \\
\midrule
Male Portrait (portrait) & Historical gentleman portrait, slight head turn, neutral expression, subtle eye movement, formal pose, period lighting & The nobleman's gaze carries quiet authority hardened by duty, a flicker of resolve crossing his weathered features as the weight of responsibility settles behind composed dignity \\
\midrule
Burghley House (architecture) & Grand manor house exterior, clouds moving overhead, trees swaying, warm afternoon sunlight, shadows shifting across stone facade & Majestic grandeur radiates from the ancient estate, centuries of ambition and legacy breathing through its towering walls as golden light bestows a reverent glow upon weathered stone \\
\midrule
Cherry Blossom (nature) & Japanese path lined with cherry trees, petals falling slowly, light filtering through branches, gentle breeze moving leaves & Ephemeral beauty drifts through the air as cherry blossoms surrender to the gentle breeze, bittersweet impermanence painting the path in soft pink as fleeting joy settles like a whisper \\
\bottomrule
\end{tabular}
\caption{Complete STOCK/NEURO prompt pairs for the diverse image generalization benchmark. STOCK prompts describe physical motion; NEURO prompts describe emotional affect grounded in the scene.}
\label{tab:diverse_prompts}
\end{table*}

\subsection{Per-Seed LPIPS Results}

\begin{table}[h]
\centering
\small
\begin{tabular}{lccc}
\toprule
\textbf{Image} & \textbf{Seed 1} & \textbf{Seed 2} & \textbf{Seed 3} \\
\midrule
mountain\_lake & 0.0391 & 0.0262 & 0.0300 \\
medieval\_market & 0.0444 & 0.0351 & 0.0320 \\
male\_portrait & 0.0148 & --- & --- \\
burghley\_house & --- & --- & --- \\
cherry\_blossom & --- & --- & --- \\
\bottomrule
\end{tabular}
\caption{Per-seed LPIPS (AlexNet) for diverse image pairs. Seeds: 42424242, 42425242, 42426242.}
\label{tab:diverse_perseed}
\end{table}

\textbf{Experimental parameters}: STOCK condition uses 30 steps, guidance 7.5, max\_shift 2.05, base\_shift 0.95. NEURO condition uses 24 steps, guidance 8.0, max\_shift 2.10, base\_shift 0.98. All renders at 512$\times$320 resolution, 49 frames, 24~fps, using LTX-2 19B (fp8) with Gemma~3 12B (fp4) text encoder. Source images span landscapes, architecture, portraits, and nature---no Victorian portraits are included.

{\small
\bibliographystyle{ieee_fullname}
\makeatletter\def\newblock{\hskip .11em plus .33em minus .07em}\makeatother
\begin{thebibliography}{1}

\bibitem{zhang2018}
R.~Zhang, P.~Isola, A.~A. Efros, et al.
\newblock The unreasonable effectiveness of deep features as a perceptual metric.
\newblock {\em CVPR}, 2018.

\end{thebibliography}
}

\end{document}
